\chapter{Remaining FYP Plan and FYP-2 Horizon}
\label{chap:fyp2}

\section{Remaining Project Scope}
With the fundamental gameplay framework, pure-GDScript neural inference, 5-floor cold-start calibration, and automated headless CI pipeline successfully established and verified at Mid-1, development transitions into subsequent iterations across the remainder of FYP-1 and FYP-2.

The remaining scope is partitioned into two major phases:
\begin{enumerate}
    \item \textbf{Phase 1: FYP-1 Final Horizon (Weeks 11--14 / Gate 3 Exit):} Delivering the complete 5-archetype combat ecosystem, the rule-based BSP generator baseline, SQLite telemetry persistence, the empirical supervised Judge model ($\text{MAE} \le 0.15$), and live polynomial TreeSHAP explainability in Godot.
    \item \textbf{Phase 2: FYP-2 Horizon (Weeks 15--32 / Gates 4 \& 5):} The in-run equipment shop economy, Judge-model re-validation, system integration and performance benchmarking, and a formal 10-participant human User Acceptance Testing (UAT) study.
\end{enumerate}

\section{Remaining Modules and Features}

\subsection{Full 5-Archetype Combat Ecosystem (Final-1)}
Expanding the enemy roster from the current 2 archetypes to 5 distinct behavioral classes:
\begin{itemize}
    \item \textbf{Melee Crawler (Completed):} Rapid pursuit, contact bite.
    \item \textbf{Ranged Pyromancer (Completed):} Standoff distance, telegraphed fireball projectile.
    \item \textbf{Heavy Brute (New):} High hit-points ($120$\,HP), slow speed ($70$\,px/s), area-of-effect ground slam shockwave.
    \item \textbf{Agile Skirmisher (New):} Flanking movement, hit-and-run sword thrusts.
    \item \textbf{Swarmers (New):} Low-health, high-count pack adversaries that flank the player.
\end{itemize}

\subsection{In-Run Equipment Shop Economy (Mid-2)}
Players collect combat-earned currency (\textit{Forge Embers}) by defeating enemies and clearing rooms. At periodic shop rooms, players can purchase:
\begin{itemize}
    \item Weapons: Increasing damage output and modifying attack swing arcs.
    \item Armor: Reducing incoming damage by a percentage.
    \item Boots: Increasing base movement speed and shortening dash cooldown.
\end{itemize}
The PPO generator will passively detect these player equipment power shifts through telemetry and dynamically scale encounter difficulty to match the augmented capability.

\subsection{Rule-Based BSP Baseline Generator (Final-1)}
To provide a classic procedural baseline alongside the Static Roguelike Baseline, the team will implement a standard Binary Space Partitioning (BSP) generator. This enables three-way comparative evaluation:
\begin{itemize}
    \item Condition A: Static Roguelike Baseline (same generator, fixed $+8\%$ per-floor stat scaling, no PPO; Non-Adaptive).
    \item Condition B: Deterministic Rule-Based BSP Generator (Non-Adaptive).
    \item Condition C: FORGE, the Two-Stage Closed-Loop PPO Generator (Player-Conditioned Adaptive).
\end{itemize}

\subsection{Empirical Judge Model and TreeSHAP Explainability (Final-1)}
Following the committee's guidance, this module has been moved forward from FYP-2 so that it is delivered and evaluated at the FYP-1 Final evaluation. It replaces the linear difficulty proxy with an offline supervised Gradient Boosted Decision Tree (LightGBM/XGBoost) trained on an empirical corpus of 25+ human playtester sessions ($\text{MAE} \le 0.15$). The model will feed into a live polynomial TreeSHAP engine in Godot~\cite{lundberg2017unified, lundberg2020local}, generating real-time \textit{Design Insight Reports} showing feature attributions for difficulty adjustments. Playtest sessions are collected as anonymised numeric telemetry files, consistent with the zero-PII governance of Section~\ref{sec:ethics}.

\section{Remaining Iterations and Development Sprints}
The upcoming development schedule is structured into three formal iterations:
\begin{itemize}
    \item \textbf{Iteration 3 (Weeks 11--14 / FYP-1 Final):}
    \begin{itemize}
        \item Sprint 3.1 (Weeks 11--12): Implement Brute, Skirmisher, and Swarmer archetypes; distribute the playtest build and begin collecting the 25-session telemetry corpus.
        \item Sprint 3.2 (Weeks 13--14): Close the corpus; train and validate the supervised Judge regressor; port TreeSHAP attribution into GDScript and build the visual explainability inspector; implement BSP generator baseline; migrate telemetry to SQLite; complete FYP-1 Final Report.
    \end{itemize}
    \item \textbf{Iteration 4 (Weeks 15--26 / FYP-2 Mid — Release Readiness):}
    \begin{itemize}
        \item Sprint 4.1 (Weeks 15--18): Build the in-run equipment shop economy (Forge Embers, shop rooms, weapons, armor, boots); extend the Judge corpus with additional playtester sessions that include equipment.
        \item Sprint 4.2 (Weeks 19--22): Re-validate and, if needed, retrain the Judge regressor against the equipment-augmented gameplay; refine the explainability inspector and Design Insight Reports; run a pre-UAT pilot session.
        \item Sprint 4.3 (Weeks 23--26): Complete system integration testing; verify non-functional performance benchmarks.
    \end{itemize}
    \item \textbf{Iteration 5 (Weeks 27--32 / FYP-2 Final — Final Acceptance):}
    \begin{itemize}
        \item Sprint 5.1 (Weeks 27--29): Conduct 10-participant human UAT study with GEQ and PENS psychometrics.
        \item Sprint 5.2 (Weeks 30--32): Statistical data analysis (paired t-tests, ANOVA); author Final Thesis; prepare public release.
    \end{itemize}
\end{itemize}

\section{Testing, Validation and Psychometric Evaluation Plan}
In FYP-2, validation expands beyond automated tests to formal human evaluation:
\begin{itemize}
    \item \textbf{Participant Cohort:} 10 human playtesters stratified across novice, intermediate, and veteran gaming experience.
    \item \textbf{Within-Subject Study Protocol:} Each participant plays three randomized 15-minute sessions:
    \begin{enumerate}
        \item Condition A: Static Baseline (+8\% per floor).
        \item Condition B: BSP Baseline (Rule-based non-adaptive).
        \item Condition C: FORGE Adaptive PPO (Closed-loop player-conditioned).
    \end{enumerate}
    \item \textbf{Standardized Psychometric Instruments:}
    \begin{itemize}
        \item \textbf{Game Engagement Questionnaire (GEQ; compared with related instruments in~\cite{denisova2016convergence}):} Assessing Absorption, Flow, Presence, and Immersion.
        \item \textbf{Player Experience of Need Satisfaction (PENS)~\cite{tyack2020sdt}:} Assessing Perceived Competence, Autonomy, and Intuitive Control.
    \end{itemize}
    \item \textbf{Statistical Hypothesis Testing:} Paired-sample t-tests and one-way repeated-measures ANOVA to establish whether FORGE achieves statistically significant gains in player flow ($p < 0.05$).
\end{itemize}

\section{Deployment Plan}
Project FORGE will be distributed as an open-source, reproducible standalone package:
\begin{itemize}
    \item \textbf{Target Platforms:} Standalone 64-bit Windows executable (\texttt{.exe}), Linux 64-bit binary (\texttt{.x86\_64}), and WebAssembly HTML5 export.
    \item \textbf{Distribution Channels:} Public itch.io game page for player recruitment and GitHub releases for open-source code and trained policy weights.
\end{itemize}

\section{Tentative Milestone Timeline}
Table~\ref{tab:gantt_plan} presents the two-semester timeline leading to final project defense.

\begin{table}[!htbp]
\centering
\caption{Two-Semester Milestone Timeline and Delivery Horizon}
\label{tab:gantt_plan}
\small
\begin{tabular}{p{0.18\textwidth}p{0.14\textwidth}p{0.48\textwidth}p{0.12\textwidth}}
\toprule
\textbf{Milestone} & \textbf{Target Date} & \textbf{Key Deliverables \& Objectives} & \textbf{Status} \\
\midrule
Proposal Defense & Week 4 & Problem statement, POC spike, proposal document & Completed \\
Mid-1 Evaluation & Week 10 & 2 enemy types, cold start, static baseline, CI & \textbf{Completed} \\
Final-1 Evaluation & Week 14 & 5 archetypes, BSP baseline, SQLite, empirical Judge model, TreeSHAP & Planned \\
Mid-2 Evaluation & Week 26 & Equipment shop economy, Judge re-validation, system integration, performance benchmarks & Planned \\
Final-2 Evaluation & Week 32 & 10-player UAT study, GEQ/PENS validation, public build & Planned \\
\bottomrule
\end{tabular}
\end{table}

\section{Expected Final Deliverables}
At the summative FYP-2 defense, Project FORGE will submit:
\begin{enumerate}
    \item Fully playable, standalone 2D action roguelike game on Windows, Linux, and Web.
    \item Complete, documented Godot 4.x and Python Gymnasium source repositories with automated CI pipelines.
    \item Trained PPO policy weights and empirical Judge regression models.
    \item Comprehensive UAT evaluation dataset with GEQ/PENS statistical findings.
    \item Final Year Project Thesis formatted to FAST-NUCES standards.
\end{enumerate}
